\clearpage
\begin{landscape}
\phantomsection\addcontentsline{toc}{section}{{\cjkfont 晶硅 有机与钙钛矿的反偏比较}}
\begin{center}
{\large\bfseries {\cjkfont 晶硅}\allowbreak{} \allowbreak{}{\cjkfont 有机与卤化物钙钛矿太阳能电池的反偏比较}\allowbreak{}\par}
\vspace{5pt}
{\fontsize{10.5}{14}\selectfont {\cjkfont 电流上翘、强反向导电与永久损伤须区分；文献电压实例不是普适阈值。温度比较使用反向电压绝对值，并固定判据。}\allowbreak{}\par}
\end{center}
\begingroup\fontsize{10.5}{14}\selectfont\setstretch{1}\renewcommand{\arraystretch}{1.05}
\setlength{\tabcolsep}{6pt}
\begin{tabular}{@{}>{\raggedright\arraybackslash}p{22mm}>{\raggedright\arraybackslash}p{66mm}>{\raggedright\arraybackslash}p{66mm}>{\raggedright\arraybackslash}p{66mm}@{}}
\toprule
\textbf{{\cjkfont 比较项}\allowbreak{}} & \textbf{{\cjkfont 晶硅}\allowbreak{}} & \textbf{{\cjkfont 有机}\allowbreak{} \allowbreak{}{\cjkfont 无序体异质结为主}\allowbreak{}} & \textbf{{\cjkfont 卤化物钙钛矿}\allowbreak{}}\\\midrule
{\cjkfont 主要机制}\allowbreak{} & {\cjkfont 雪崩倍增；高掺杂窄结可有带间隧穿；缺陷和局部高场可提前导通}\allowbreak{} [C1,C2] & {\cjkfont 候选为接触注入}\allowbreak{}/\allowbreak{}{\cjkfont 隧穿、陷阱辅助或其他场增强产生、局部弱区导电。本项目机制尚未确定}\allowbreak{} [C3] & {\cjkfont 离子再分布改变界面场；叠加注入}\allowbreak{}/\allowbreak{}{\cjkfont 隧穿、电化学反应和局部热效应}\allowbreak{} [C5,C6]\\[5pt]
{\cjkfont 输运特点}\allowbreak{} & {\cjkfont 以能带输运为主；雪崩需满足碰撞电离条件，不能仅凭“自由程长”判断}\allowbreak{} & {\cjkfont 常用局域态跳跃描述，受无序、占据及形貌连通性控制；跳距不等于能带自由程}\allowbreak{} [C4] & {\cjkfont 电子}\allowbreak{}/\allowbreak{}{\cjkfont 空穴输运与可移动离子并存；离子改变场、势垒和化学状态}\allowbreak{}\\[5pt]
{\cjkfont 触发场和电压}\allowbreak{} & {\cjkfont 取决于掺杂、耗尽层、结曲率与缺陷；局部提前导通不等于理想均匀结击穿}\allowbreak{} & 10 V \allowbreak{}{\cjkfont 若全落在}\allowbreak{} 100 nm \allowbreak{}{\cjkfont 层上，平均场为}\allowbreak{} 1 MV/cm\nobreak{}{\cjkfont ，非普适阈值。优化}\allowbreak{} OSC \allowbreak{}{\cjkfont 有不可逆击穿}\allowbreak{} \ensuremath{\vert}V\ensuremath{\vert}\ensuremath{>}35 V \allowbreak{}{\cjkfont 的报道}\allowbreak{} [C3] & {\cjkfont 特定早期器件为}\allowbreak{} −1 \allowbreak{}{\cjkfont 至}\allowbreak{} −4 V [C5]\nobreak{}{\cjkfont ；优化阻挡层已有}\allowbreak{} \ensuremath{\vert}Vbd\ensuremath{\vert}\ensuremath{>}20 V [C6]\nobreak{}{\cjkfont ，不宜给统一范围}\allowbreak{}\\[5pt]
{\cjkfont 温度依赖}\allowbreak{} & {\cjkfont 理想雪崩主导时}\allowbreak{} \ensuremath{\vert}Vbd\ensuremath{\vert} \allowbreak{}{\cjkfont 通常随温度升高而增大；}\allowbreak{}Zener \allowbreak{}{\cjkfont 主导者可减小，混合机制另判}\allowbreak{} [C1] & {\cjkfont 发射}\allowbreak{}/\allowbreak{}{\cjkfont 注入与隧穿响应不同；固定电流阈值和形状膝点不可混用。本项目尚无变温定论}\allowbreak{} & {\cjkfont 离子、反应和电子过程共同影响；表观阈值也随测试时长及预处理而变}\allowbreak{}\\[5pt]
{\cjkfont 时间和历史依赖}\allowbreak{} & {\cjkfont 电子响应可快，但自热、缺陷演化和应力损伤会产生历史依赖}\allowbreak{} [C2] & {\cjkfont 本项目：扫描快慢及是否脉冲对膝点影响不明显；脉宽等待补齐，不能泛化为所有器件}\allowbreak{} & {\cjkfont 常有明显偏压历史、保持时间和扫描程序影响；程度依结构及条件而变}\allowbreak{} [C5,C6]\\[5pt]
{\cjkfont 损伤与可逆性}\allowbreak{} & {\cjkfont 局部电流可形成热点并永久损伤；反向导电开始不等于已损坏}\allowbreak{} & {\cjkfont 本项目：轻微越过膝点后，正向}\allowbreak{} J–V \allowbreak{}{\cjkfont 基本不变，反扫可重现。更强应力的永久损伤机制另查}\allowbreak{} & {\cjkfont 可有可恢复变化或永久短路、金属丝、腐蚀与分解；不保证一律“先软后硬”}\allowbreak{} [C5,C6]\\[5pt]
{\cjkfont 常见缓解方向}\allowbreak{} & {\cjkfont 旁路保护、控制缺陷和局部场集中、改善散热并验证耐受性}\allowbreak{} & {\cjkfont 改善膜厚}\allowbreak{}/\allowbreak{}{\cjkfont 形貌均匀性、接触选择性及阻挡层，配合限流；措施需匹配机制}\allowbreak{} & {\cjkfont 离子}\allowbreak{}/\allowbreak{}{\cjkfont 金属阻挡层、稳定接触与传输层、减少局部缺陷电流并配合模块保护}\allowbreak{} [C6]\\[5pt]
\bottomrule\end{tabular}\par\endgroup
\noindent\vspace{6pt}{\fontsize{9.5}{12}\selectfont {\cjkfont 本项目观察来自用户陈述，尚未逐项核对原始曲线和脉冲波形。实际局部电场需自洽求解；银迁移不是本器件已证实的损伤机制。资料索引见下一页。}\allowbreak{}\par}\end{landscape}\clearpage
\subsection*{{\cjkfont 比较表的来源与使用边界}\allowbreak{}}
{\cjkfont 这张表用于建立候选机制和测量尺度的比较，不用于仅凭端电压、温度趋势或曲线形状给器件定性。特别是本项目已报告轻微越过膝点具有可重复性，应先解释可逆的电流上翘，再单独研究更强应力引发的损伤。}\allowbreak{}\par\medskip
\noindent [C1] M. Singh Tyagi (1968). Zener and avalanche breakdown in silicon alloyed p–n junctions—II. Solid-State Electronics 11, 117–128. DOI: \href{https://doi.org/10.1016/0038-1101(68)90142-1}{\nolinkurl{10.1016/0038-1101(68)90142-1}}\par\smallskip
\noindent [C2] Influence of surface texture on the defect-induced breakdown behavior of multicrystalline silicon solar cells. Progress in Photovoltaics (2013). DOI: \href{https://doi.org/10.1002/pip.1226}{\nolinkurl{10.1002/pip.1226}}\par\smallskip
\noindent [C3] J. Huang et al. (2026). Perovskite–organic tandem solar cells with superior reverse-bias stability. Nature Materials 25, 1419–1428. DOI: \href{https://doi.org/10.1038/s41563-026-02541-6}{\nolinkurl{10.1038/s41563-026-02541-6}}\par\smallskip
\noindent [C4] W. F. Pasveer et al. (2005). Unified Description of Charge-Carrier Mobilities in Disordered Semiconducting Polymers. Physical Review Letters 94, 206601. DOI: \href{https://doi.org/10.1103/PhysRevLett.94.206601}{\nolinkurl{10.1103/PhysRevLett.94.206601}}\par\smallskip
\noindent [C5] A. R. Bowring, L. Bertoluzzi, B. C. O'Regan, M. D. McGehee (2017). Reverse Bias Behavior of Halide Perovskite Solar Cells. Advanced Energy Materials, 1702365. DOI: \href{https://doi.org/10.1002/aenm.201702365}{\nolinkurl{10.1002/aenm.201702365}}\par\smallskip
\noindent [C6] Barrier reinforcement for enhanced perovskite solar cell stability under reverse bias. Nature Energy 9, 1264–1274 (2024). DOI: \href{https://doi.org/10.1038/s41560-024-01579-7}{\nolinkurl{10.1038/s41560-024-01579-7}}\par\smallskip
\medskip{\small {\cjkfont 证据范围：}\allowbreak{}C3 \allowbreak{}{\cjkfont 数值取自作者机构公布的原始摘要，仅指该研究体系；}\allowbreak{}C5 \allowbreak{}{\cjkfont 与}\allowbreak{} C6 \allowbreak{}{\cjkfont 为不同结构和测试条件。}\allowbreak{}C2 \allowbreak{}{\cjkfont 的损伤补充依据为}\allowbreak{} Thermal and electrical investigation of the reverse bias degradation of silicon solar cells\allowbreak{}{\cjkfont （}\nobreak{}Microelectronics Reliability\nobreak{}{\cjkfont ，}\allowbreak{}2013\nobreak{}{\cjkfont ）。}\allowbreak{}\par}\clearpage
